
%BeginMC
\question

\begin{choices}
\correctchoice  %EndChoice
\choice  %EndChoice
\choice  %EndChoice
\choice  %EndChoice
\choice  %EndChoice
\end{choices}
%EndMC

%%%%%%%%%%%% Questions begin here. Copy and paste from template above.

%BeginMC
\question
While observing a population of lizards during one week, you notice that the number of adults has increased and is higher at the end of the week than at the start. The number of juveniles is unchanged. The most likely explanation for such an observation would be 

\begin{choices}
\correctchoice increased immigration.  %EndChoice
\choice decreased death rate. %EndChoice
\choice increased emigration.  %EndChoice
\choice decreased immigration. %EndChoice
\choice decreased emigration.  %EndChoice
\choice increased birth rate.  %EndChoice
\end{choices}
%EndMC
	
%BeginMC
\question
A secondary consumer has high abundance. In a trophic cascade, which of the following is probably \emph{not} true.
\begin{choices}
\correctchoice The primary producer will have low abundance. %EndChoice
\choice The tertiary consumer will have low abundance. %EndChoice
\choice The primary consumer will have low abundance. %EndChoice
\choice The different trophic levels will have alternating high and low abundances. %EndChoice
\choice Actually, all of the above are true. %EndChoice%dontshuffle
\end{choices}
%EndMC



%BeginMC
\question
Which of the following statements is consistent with the principle of competitive exclusion?  

\begin{choices}
\correctchoice Even a slight reproductive advantage will eventually lead to the elimination of the less well adapted of two competing species. %EndChoice
\choice Evolution tends to increase competition between related species.   %EndChoice
\choice Bird species generally do not compete for nesting sites.  %EndChoice
\choice The density of one competing species will have a positive impact on the population growth of the other competing species. %EndChoice
\choice Two species with the same fundamental niche will exclude other competing species.   %EndChoice
\end{choices}
%EndMC


%BeginMC
\question
The attribute that most influence population density and rates of immigration and emigration among habitat patches in a metapopulation are

\begin{choices}
\correctchoice patch size, patch isolation and habitat quality. %EndChoice
\choice patch size, patch isolation and species diversity. %EndChoice
\choice patch isolation, habitat quality, and size of source patches. %EndChoice
\choice patch size, habitat quality, and isolation of sink patches. %EndChoice
\choice habitat quality, size of sink patches and isolation of source patches. %EndChoice
\end{choices}
%EndMC

	
%BeginMC
\question
Chimpanzees and humans may be as much as 98.5\% genetically identical.  A likely reason for the large phenotypic difference between chimps and humans is

\begin{choices}
\correctchoice a difference in the developmental expression and timing of toolkit genes. %EndChoice
\choice the similar genotype interacting with the different physical environments occupied by each species. %EndChoice
\choice a large difference in the types of genes found in each species. %EndChoice
\choice paedomorphic development of most human traits. %EndChoice
\choice allometric development of most human traits. %EndChoice
\choice heterochronic development of most chimpanzee traits. %EndChoice
\end{choices}
%EndMC

%BeginMC
\question
A change in the body location where two \textit{Hox} genes are expressed caused swimming appendages to evolve into feeding appendages in crabs and other crustaceans. This type of genetic change is an example of

\begin{choices}
\correctchoice a change in a toolkit gene that altered the arrangement of parts in the body plan. %EndChoice
\choice heterochrony that caused a difference in developmental timing. %EndChoice
\choice allometry that caused the swimming appendages to grow more slowly so they could be used for feeding. %EndChoice
\choice paedomorphosis that caused the adult to retain the juvenile feeding appendage. %EndChoice
\choice an adaptation to living in an quiet aquatic environment where fewer swimming appendages are needed. %EndChoice
\end{choices}
%EndMC


%BeginMC
\question
The basic model of metapopulations views the number of suitable habitat patches occupied at a given point in time as a balance between

\begin{choices}
\correctchoice recolonization and local extinction. %EndChoice
\choice dispersal and birth rate. %EndChoice
\choice birth rate and local extinction. %EndChoice
\choice recolonization and degree of isolation. %EndChoice
\choice resource quality and degree of isolation. %EndChoice
\end{choices}
%EndMC


%BeginMC
\question
Symbiosis is a longer-term ecological interaction

\begin{choices}
\correctchoice between two or more species. %EndChoice
\choice where one species benefits and another species is harmed %EndChoice
\choice where both species benefit. %EndChoice
\choice  where one species benefits and the other species neither benefits nor is harmed. %EndChoice
\choice where two or more species partition resources. %EndChoice
\end{choices}
%EndMC


%BeginMC
\question
Competitive exclusion means that

\begin{choices}
\correctchoice one species will always out complete another species that occupies the same niche.  %EndChoice
\choice two species will partition resources to eliminate competition. %EndChoice
\choice one individual will always out complete another individual in a population. %EndChoice
\choice one population will always exclude another another population of the same species in a community. %EndChoice
\choice a species will have a smaller realized niche than a fundamental niche. %EndChoice
\end{choices}
%EndMC

%BeginMC
\question
A field containing 2000 kg of plant material can support about \rule{1in}{0.4pt} of tertiary consumers.

\begin{choices}
\correctchoice 2 kg %EndChoice
\choice 20 kg %EndChoice
\choice 10 kg %EndChoice
\choice 10\% %EndChoice
\choice 200 kg %EndChoice
\choice 90\% %EndChoice
\end{choices}
%EndMC


%BeginMC
\question
The goal of hierarchical classification is to

\begin{choices}
\choice provide an easy to memorize arrangement of life on earth.
\choice provide a testable hypothesis about organisms with similar structures.
\choice use a phylogenetic tree to accurately name taxa.
\correctchoice create monophyletic groups that represent evolutionary history.
\choice represent all life on earth with only eight categories: domain, kingdom, phylum, class, order, family, genus, species.
\end{choices}
%EndMC

%BeginMC
\question
In a mature forest of oak, maple, and hickory trees, a disease causes a reduction in the number of acorns produced by oak trees. Which of the following would \emph{least} likely be a \emph{direct} result of this? 

\begin{choices}
\correctchoice There might be a decrease in the number of maple seeds as the disease spreads to other trees in the forest. %EndChoice
\choice There might be fewer squirrels because they feed on acorns produced by oak trees. %EndChoice
\choice  There might be fewer mice and seed-eating birds because squirrels that normally eat acorns would eat more seeds and compete with the mice and birds.%EndChoice
\choice There might be an increase in the number of hickory trees because the competition between hickory nuts and acorns for germination sites would be reduced or eliminated. %EndChoice
\choice There might be fewer owls because they feed on baby squirrels, mice, and young seed-eating birds, whose populations would be reduced. %EndChoice
\end{choices}
%EndMC


%BeginMC
\question
A small population of white-footed mice has the same intrinsic rate of increase (\textit{r}) as a large population.  If everything else is equal, 

\begin{choices}
\correctchoice the large population will add more individuals per unit time.  %EndChoice
\choice the small population will add more individuals per unit time. %EndChoice
\choice the two populations will add equal numbers of individuals per unit time.%EndChoice
\choice the growth trajectories of the two populations will proceed in opposite directions. %EndChoice
\choice the number of individuals added will depend on the value of N. %EndChoice
\end{choices}
%EndMC


%BeginMC
\question
Which of the following areas of study focuses on the exchange of energy and nutrients among individuals of different species?

\begin{choices}
\correctchoice Ecosystem ecology. %EndChoice
\choice Population Ecology %EndChoice
\choice Evolutionary Ecology %EndChoice
\choice Community Ecology %EndChoice
\choice Behavioral Ecology %EndChoice
\end{choices}
%EndMC


%BeginMC
\question
A patch in a population has the following four values. 

\qquad\begin{tabular}{@{}ll@{}}
$b = 0.04$ & $i = 0.06$\\
$d = 0.08$ & $e = 0.04$\\
\end{tabular}

Based only on these values (ignoring habitat quality, patch size, etc.), the patch is most likely to be a 

\begin{choices}
\correctchoice  sink.%EndChoice
\choice  source. %EndChoice
\end{choices}
%EndMC



%BeginMC
\question
Assume that you performed an \textsc{anova}. The final probability was $p = 0.035$. Based on this information, which of the following is true?

\begin{choices}
\correctchoice You would reject the null hypothesis that the samples are from the same statistical population. %EndChoice
\choice You would accept the null hypothesis that the samples are from the same statistical population. %EndChoice
\choice You would reject the alternative hypothesis that the samples are from the same statistical population. %EndChoice
\choice You would reject the alternative hypothesis that the samples are from different statistical populations.
 %EndChoice
\choice You would accept the null hypothesis that the samples are from different statistical populations. %EndChoice
\choice You would accept the alternative hypothesis that the samples are from the same statistical population. %EndChoice
\end{choices}
%EndMC


%BeginMC
\question
Two species that each share an immediate common ancestor are called

\begin{choices}
\correctchoice sister species. %EndChoice
\choice descendant species. %EndChoice
\choice taxa. %EndChoice
\choice branch species. %EndChoice
\choice rooted species. %EndChoice
\end{choices}
%EndMC

%BeginMC
\question
The science of naming a new species is called

\begin{choices}
\correctchoice  taxonomy. %EndChoice
\choice binomial nomenclature. %EndChoice
\choice classification. %EndChoice
\choice systematics. %EndChoice
\choice phylogeny. %EndChoice
\end{choices}
%EndMC



%BeginMC
\question
Which of the following is the correct and complete arrangement of the hierarchical classification. 
\begin{choices}
\correctchoice  Domain, Kingdom, Phylum, Class, Order, Family, Genus, Species.%EndChoice
\choice  Domain, Kingdom, Phylum, Family, Class, Species, Genus.%EndChoice
\choice  Domain, Kingdom, Phylum, Order, Class, Family, Genus, Species.%EndChoice
\choice  Domain, Kingdom, Phylum, Order, Class, Genus, Family, Species.%EndChoice
\choice  Domain, Phylum, Family, Class, Order, Genus, Species.%EndChoice
\end{choices}
%EndMC

%BeginMC
\question
Which block in the figure below shows the greatest species diversity? (Different shades, including white, represent different species. The small squares indicate relative abundance of individuals present.)

\begin{multicols}{3}
\begin{choices}
\correctchoice \includegraphics{diversity_mcE} %EndChoice
\choice \includegraphics{diversity_mcA} %EndChoice
\choice \includegraphics{diversity_mcB} %EndChoice
\choice \includegraphics{diversity_mcC} %EndChoice
\choice \includegraphics{diversity_mcD} %EndChoice
\end{choices}
\end{multicols}
%EndMC

%BeginMC
\question
Competition for limited resources may influence gene flow by causing

\begin{choices}
\correctchoice some individuals to migrate to another population.%EndChoice
\choice  some individuals to partition resources, reducing intraspecific competition.%EndChoice
\choice  an increased death rate, thereby reducing gene flow.%EndChoice
\choice  some individuals to move to a different source patch in the same population.%EndChoice
\choice  some individuals to evolve a new trait to be better competitors.%EndChoice
\end{choices}
%EndMC

%BeginMC
\question
White-breasted nuthatches and brown creepers both eat insects that hide in the furrows of bark in hardwood trees. The brown creeper searches for insects by hunting from the bottom of the tree trunk toward the top, whereas the white-breasted nuthatch searches from the top of the trunk down. Their feeding areas rarely overlap on the tree trunk. These hunting behaviors best illustrate which of the following ecological concepts?

\begin{choices}
\choice competitive exclusion%EndChoice
\correctchoice resource partitioning%EndChoice
\choice intraspecific competition%EndChoice
\choice keystone species%EndChoice
\choice interspecific competition%EndChoice
\end{choices}
%EndMC


%BeginMC
\question
Given the following values and assuming exponential population growth, what will be the total population size after one generation has elapsed? Round to nearest whole individual.\vspace{1ex}

\begin{minipage}{1em}
\end{minipage}%
\hspace{2em}%
\begin{minipage}{1.5in}
$b = 0.09$ \\
$d = 0.12$ \\
$N = 18,657$\\
$K = 15,000$\\
\end{minipage}

\begin{choices}
\choice $15,000$%EndChoice
\choice $560$%EndChoice
\correctchoice $18,097$%EndChoice
\choice $19,217$%EndChoice
\choice $3,657$%EndChoice
\end{choices}
%EndMC


%BeginMC
\question
In the experiment shown in the figure below, \textit{Balanus} was removed from the habitat shown on the left. Which of the following statements is a valid conclusion of this experiment?\\
\includegraphics[width=0.9\textwidth]{balanus.png}

\begin{choices}
\choice \textit{Balanus} can survive only in the lower intertidal zone because it is unable to resist dehydration.%EndChoice
\choice Balanus is inferior to \textit{Chthamalus} in competing for space on rocks lower in the intertidal zone. %EndChoice
\correctchoice The removal of \textit{Balanus} shows that the realized niche of \textit{Chthamalus} is smaller than its fundamental niche.%EndChoice
\choice These two species of barnacle do not show competitive exclusion.%EndChoice
\choice If \textit{Chthamalus} were removed, \textit{Balanus}'s fundamental niche would become larger.%EndChoice
\end{choices}
%EndMC


%BeginMC
\question
Based on the hypothetical food web below, one letter represents an organism that could be both a secondary and tertiary consumer? The arrows point in the direction of energy transfer from lower to higher trophic levels. More than one answer is possible. Circle only one answer.

\begin{multicols}{3} 
\begin{choices}
\correctchoice A%EndChoice
\choice B%EndChoice
\choice C%EndChoice
\choice D%EndChoice
\correctchoice E%EndChoice
\end{choices}
\columnbreak
\includegraphics{foodweb.png}
\columnbreak
\end{multicols}
\vspace{-1\baselineskip}
%EndMC

%% Same as above question but using trophic.
%BeginMC
\question
Based on the hypothetical food web below, one letter represents an organism that could be both a secondary and tertiary consumer? The arrows point from lower to higher trophic levels. More than one answer is possible. Circle only one answer.

\begin{multicols}{3} 
\begin{choices}
\correctchoice A%EndChoice
\choice B%EndChoice
\choice C%EndChoice
\choice D%EndChoice
\correctchoice E%EndChoice
\end{choices}
\columnbreak
\includegraphics{foodweb.png}
\columnbreak
\end{multicols}
\vspace{-1\baselineskip}
%EndMC
 
%BeginMC
\question
Based on the hypothetical food web below, which letter represents an organism that is most likely a detritivore?  The arrows point to organisms receiving energy.

\begin{multicols}{3} 
\begin{choices}
\correctchoice A%EndChoice
\choice B%EndChoice
\choice C%EndChoice
\choice D%EndChoice
\choice E%EndChoice
\end{choices}
\columnbreak
\includegraphics{foodweb.png}
\columnbreak
\end{multicols}
\vspace{-1\baselineskip}
%EndMC
 
 %BeginMC
\question
Based on the hypothetical food web below, which letter represents an organism that is most likely a primary producer?  The arrows point to organisms receiving energy.

\begin{multicols}{3} 
\begin{choices}
\choice A%EndChoice
\correctchoice B%EndChoice
\choice C%EndChoice
\choice D%EndChoice
\choice E%EndChoice
\end{choices}
\columnbreak
\includegraphics{foodweb.png}
\columnbreak
\end{multicols}
\vspace{-1\baselineskip}
%EndMC
 


%BeginMC
\question
Which of the following statements about ecology is \emph{incorrect}? 

\begin{choices}
\choice Ecologists may study populations and communities of organisms. %EndChoice
\correctchoice Ecology is a discipline that is independent from natural selection and evolutionary history. %EndChoice
\choice Ecology spans increasingly comprehensive levels of organization, from individuals to ecosystems. %EndChoice
\choice Ecology includes the study of the interactions between biotic and abiotic aspects of the environment. %EndChoice
\choice Ecologists may study the behavioral interactions between two species.%EndChoice
\end{choices}
%EndMC

%BeginMC
\question
In ecosystems, why is the term cycling used to describe the transfer of chemical elements and nutrients, whereas the term flow is used for the transfer of energy?

\begin{choices}
\correctchoice Chemical elements are recycled within ecosystems, but energy flows through and out of ecosystems.%EndChoice
\choice Both chemical elements and energy are recycled but can flow to other ecosystems.%EndChoice
\choice Chemical elements are cycled into ecosystems from other ecosystems, but energy constantly flows within the ecosystem.%EndChoice
\choice Both chemical elements and energy flow within an ecosystem but can cycle to other ecosystems.%EndChoice
\choice Energy flows from one ecosystem to other ecosystems, but chemical elements constantly cycle within an ecosystem.%EndChoice
\end{choices}
%EndMC

%BeginMC
\question
Consider the food chain grass $\to$ grasshopper $\to$ mouse $\to$ snake $\to$ hawk. How much of the chemical energy fixed by photosynthesis of the grass (100\%) is available to the hawk?

\begin{choices}
\correctchoice 0.01\%%EndChoice
\choice 0.1\%%EndChoice
\choice 1\%%EndChoice
\choice 10\%%EndChoice
\choice 60\%%EndChoice
\end{choices}
%EndMC

%BeginMC
\question
The goal of hierarchical classification is to

\begin{choices}
\choice provide an easy to memorize arrangement of life on earth.%EndChoice
\choice provide a testable hypothesis about organisms with similar structures.%EndChoice
\choice use a phylogenetic tree to accurately name taxa.%EndChoice
\correctchoice create monophyletic groups that represent evolutionary history.%EndChoice
\choice represent all life on earth with only eight categories: domain, kingdom, phylum, class, order, family, genus, species.%EndChoice
\end{choices}
%EndMC

%BeginMC
\question
A genetic change that caused a certain \textit{Hox} gene to be expressed along the tip of a vertebrate limb bud instead of farther back helped make possible the  evolution of the tetrapod limb. This type of change is an example of 
\begin{choices}
\choice the influence of environment on development.%EndChoice
\choice paedomorphosis.%EndChoice
\correctchoice a change in a toolkit gene that altered the organization of body parts.%EndChoice
\choice heterochrony.%EndChoice
\choice allometry.%EndChoice
\end{choices}
%EndMC

%BeginMC
\question
Which of the following is most likely to contribute to density-dependent regulation of populations?

\begin{choices}
\choice the removal of toxic waste by decomposers%EndChoice
\correctchoice intraspecific competition for nutrients%EndChoice
\choice earthquakes%EndChoice
\choice floods%EndChoice
\choice fires%EndChoice
\end{choices}
%EndMC

%BeginMC
\question
Often the growth cycle of one population has an effect on the cycle of another. As moose populations increase, for example, wolf populations also increase. If you consider the logistic equation for the wolf population, which of the factors accounts for the effect of the moose population on the size of the wolf population?

\begin{choices}
\choice $r$%EndChoice
\choice $N$%EndChoice
\choice $rN$%EndChoice
\correctchoice $K$%EndChoice
\choice $\Delta t$%EndChoice
\choice $\Delta N$%EndChoice
\end{choices}
%EndMC


%BeginMC
\question
Naturally occurring fires cause occasional disturbance to an ecosystem. If humans stop these fires as soon as they are started, then

\begin{choices}
\choice species diversity in the ecosystem will tend to increase.%EndChoice
\correctchoice species diversity in the ecosystem will tend to decrease.%EndChoice
\choice species diversity will tend to remain unchanged, but the types of species present may change.%EndChoice
\choice the diversity of threatened and endangered species will be protected.%EndChoice
\end{choices}
%EndMC

%BeginMC
\question
Fire suppression by humans

\begin{choices}
\choice will always result in an increase in species diversity in a given biome.%EndChoice
\correctchoice can change the species composition within biological communities.%EndChoice
\choice will result ultimately in sustainable production of increased amounts of forest products for human use.%EndChoice
\choice is necessary for the protection of threatened and endangered forest species.%EndChoice
\choice is a management goal of conservation biologists to maintain the healthy condition of forest communities.%EndChoice
\end{choices}
%EndMC


%BeginMC
\question
The biological species concept emphasizes

\begin{choices}
\choice the formation of cryptic species.%EndChoice
\choice gene flow between species.%EndChoice
\choice reproductive isolation between populations.%EndChoice
\choice reproductive isolation between species.%EndChoice
\choice the formation of morphologically distinct species.%EndChoice
\end{choices}
%EndMC

